%==============================================================================
\section{How the package is built}\label{sec:architecture}
%==============================================================================

FastHydrology is plain Julia on plain arrays: no GPU backend, no preferences, no macros of its own. Its structure
rests on four ideas: one grid interface every model is written against; a state shared by all models and a
model object that holds everything model-specific; every option a type, so each choice is made once by dispatch;
and optional dependencies as package extensions.

\subsection{Abstract types and what each concrete type supplies}\label{arch:types}
\begin{center}\small
\begin{tabular}{>{\raggedright\arraybackslash}p{3.6cm}>{\raggedright\arraybackslash}p{4.3cm}>{\raggedright\arraybackslash}p{7.2cm}}
\toprule
abstract type & concrete types & contract \\
\midrule
\code{Abstract\allowbreak{}Hydro\allowbreak{}Grid} & \code{ArrayHydroGrid\{T\}} & fields \code{Nx}, \code{Ny}, \code{dx}, \code{dy}; \code{alloc\_field}; optionally the field operations of \file{common/\allowbreak{}grid.jl} \\
\code{Abstract\allowbreak{}Hydro\allowbreak{}State} & \code{HydroState\{A\}} & any fields; checkpoints use \code{fieldnames} \\
\code{Abstract\allowbreak{}Hydro\allowbreak{}Model} & \code{Kazmierczak\allowbreak{}Hydro\allowbreak{}Model}, \code{HABHydroModel}, \code{Shakti\allowbreak{}Hydro\allowbreak{}Model} & \code{update\_steady\_state!(model, grid, state)} (steady models) or \code{run!}/\code{step!} on a \code{TimeSimulation}; optionally \code{N\_responds\_to\_ub} \\
\code{Abstract\allowbreak{}Simulation} & \code{Steady\allowbreak{}State\allowbreak{}Simulation}, \code{TimeSimulation} & \code{run!} \\
\code{Abstract\allowbreak{}Output\allowbreak{}Writer} & \code{Net\allowbreak{}CDFOutput\allowbreak{}Writer} & \code{write\_output!}, \code{close\_output!} \\
\bottomrule
\end{tabular}
\end{center}

\paragraph{The grid interface.} A grid exposes \code{Nx}, \code{Ny}, \code{dx}, \code{dy} as fields, read directly by the
physics code (\code{grid.dx}); the float type is \code{typeof(grid.dx)}. The only function a new grid must implement is
\code{alloc\_field}, the one place that knows the field type. The others have plain-array defaults in
\file{common/grid.jl}: \code{fill\_halo!} (a no-op), \code{convolve!} (ImageFiltering's \code{imfilter!}),
\code{minus\_gradient\_x!/y!}, \code{masked\_mean}, \code{masked\_max\_abs}, \code{masked\_max\_abs\_diff},
\code{overwrite\_where!}. \code{ArrayHydroGrid} overrides only \code{alloc\_field} and \code{convolve!} (with the cached FFT).
The interface is the remnant of v1, where an Oceananigans-backed grid carried halo storage; that grid was removed in
v2.0.0 (commit \code{56f381e}), and the calls to \code{fill\_halo!} that remain in the physics are no-ops.

\paragraph{State versus model.} \code{HydroState} holds only what every model shares: inputs \code{mask}, \code{h}, \code{b}
and outputs \code{N}, \code{W}. Everything else (inputs such as $\kappa$, $|u_b|$, $A$, $G$, $q_T$, and every intermediate
field) lives on the model. One state type thus serves all models, and the checkpoint code needs no knowledge of any
model.

\subsection{Options as types}\label{arch:options}
Every switch of the K24 model is a value of an abstract option type, stored in the model's type parameters or fields,
and each function that depends on it has one method per value. The choice is therefore resolved by the compiler
when the model's type is known, not tested in the loops. For example, the dissipation melt:
\begin{verbatim}
abstract type AbstractDissipationMelt end
struct DissipationMeltOn  <: AbstractDissipationMelt end
struct DissipationMeltOff <: AbstractDissipationMelt end
add_dissipation_term!(model, ::DissipationMeltOff) = (model.Q_diss .= 0; nothing)
add_dissipation_term!(model, ::DissipationMeltOn)  =
    add_dissipation_term!(model, model.dissipation_discretization)
\end{verbatim}
The constructor turns the user's \code{dissipation\_melt = true} into \code{DissipationMeltOn()}. The option families:
\begin{center}\small
\begin{longtable}{>{\raggedright\arraybackslash}p{4.4cm}>{\raggedright\arraybackslash}p{5.2cm}>{\raggedright\arraybackslash}p{5.4cm}}
\toprule
abstract type & values & dispatched in \\
\midrule
\endhead
\code{Abstract\allowbreak{}Dissipation\allowbreak{}Melt} & \code{Dissipation\allowbreak{}Melt\allowbreak{}On}, \code{Dissipation\allowbreak{}Melt\allowbreak{}Off} & \code{resolve\_q!}, \code{add\_dissipation\_term!} \\
\code{Abstract\allowbreak{}Psi\allowbreak{}Out\allowbreak{}Algorithm} & \code{TapedPsiOut}, \code{RecursivePsiOut}, \code{IterativePsiOut}, \code{Topological\allowbreak{}Psi\allowbreak{}Out} & \code{route\_psi\_out!} \\
\code{Abstract\allowbreak{}Routing\allowbreak{}Scheme} & \code{Warner}, \code{GDSWarner}, \code{Quinn}, \code{Tarboton}, \code{Modified\allowbreak{}Tarboton}, \code{GDSTarboton} & \code{prepare\_routing!}, \code{cell\_weights!}, \code{n\_directions} \\
\code{Abstract\allowbreak{}Fill\allowbreak{}Algorithm} & \code{Priority\allowbreak{}Flood\allowbreak{}Fill}, \code{JacobiFill}, \code{Lowest\allowbreak{}Neighbour\allowbreak{}Fill} & \code{fill\_potential!} \\
\code{Abstract\allowbreak{}QConversion} & \code{QFrom\allowbreak{}Face\allowbreak{}Average}, \code{QFromOutflow} & \code{q\_from\_psi\_out!} \\
\code{Abstract\allowbreak{}Dissipation\allowbreak{}Discretization} & \code{FaceDissipation}, \code{Cell\allowbreak{}Centred\allowbreak{}Dissipation} & \code{add\_dissipation\_term!} \\
\code{Abstract\allowbreak{}Friction\allowbreak{}Discretization} & \code{Cell\allowbreak{}Centred\allowbreak{}Friction}, \code{Staggered\allowbreak{}Friction} & \code{add\_friction\_term!} \\
\code{Abstract\allowbreak{}Water\allowbreak{}Thickness\allowbreak{}Algorithm} & \code{Darcy\allowbreak{}Weisbach\allowbreak{}Thickness\{G\}}, \code{Laminar\allowbreak{}Thickness\{G\}}, \code{Areal\allowbreak{}Conduit\allowbreak{}Thickness} & \code{update\_W!} \\
\code{Abstract\allowbreak{}Gradient\allowbreak{}Convention} & \code{MeanGradient}, \code{LocalGradient} & \code{update\_W\_darcy\_weisbach!}, \code{update\_W\_laminar!} \\
\code{Abstract\allowbreak{}Sliding\allowbreak{}Law} (\code{Abstract\allowbreak{}Pressure\allowbreak{}Dependent\allowbreak{}Sliding\allowbreak{}Law}) & seven laws, Section~\ref{phys:sliding} & \code{update\_tau\_b!}, \code{calc\_tau\_b}, \code{resolve\_q!} \\
\code{Abstract\allowbreak{}Drainage\allowbreak{}Mode} & \code{BothDrainage}, \code{EfficientOnly}, \code{InefficientOnly} & \code{effective\_Q\_c}, \code{opening\_coefficients}, \code{N\_responds\_to\_ub} \\
\bottomrule
\end{longtable}
\end{center}
Five of these (dissipation melt, sliding law, psi-out algorithm, water thickness, drainage mode) are type parameters of
\code{KazmierczakParams}, so a model's type spells out its configuration. The fill algorithm, routing scheme, $q$
conversion and the two discretisations are stored in abstractly typed fields; their dispatch happens once per call at
a function barrier (\code{fill\_potential!}, \code{prepare\_routing!}, \code{q\_from\_psi\_out!}, \ldots), outside the hot loops.

The gradient convention is a type parameter of the thickness closure rather than a \code{Bool}, because the mean case
divides by a scalar and the local case by a field: a runtime switch would leave a small \code{Union} inside the
broadcast.

\subsection{The K24 model: parameters, workspace and \code{getproperty}}\label{arch:split}
\code{KazmierczakHydroModel} has two fields: \code{params::KazmierczakParams} (39 immutable constants and options) and
\code{workspace::KazmierczakWorkspace} (38 arrays and buffers, allocated once and updated in place). Neither struct is
\code{mutable}: nothing ever reassigns a field, only the contents of the arrays (\code{model.q .= ...}). The split keeps the
two positional constructor calls short and grouped, instead of one 80-argument call that silently breaks if a field is
inserted out of order. \code{Base.getproperty} is overloaded so that \code{model.<name>} looks first in \code{params}, then
in \code{workspace}:
\begin{verbatim}
function Base.getproperty(model::KazmierczakHydroModel, name::Symbol)
    name === :params    && return getfield(model, :params)
    name === :workspace && return getfield(model, :workspace)
    params = getfield(model, :params)
    hasfield(typeof(params), name) && return getfield(params, name)
    return getfield(getfield(model, :workspace), name)
end
\end{verbatim}
So \code{model.rho\_w} and \code{model.q} both work, and the physics code was not changed by the split. \code{hasfield} on a
concrete type is resolved at compile time, so the forwarding costs nothing. \code{propertynames} lists both sets.
Appendix~\ref{app:fields} lists every field.

\subsection{Float type}\label{arch:float}
The float type is the grid's: \code{ArrayHydroGrid(...; T = Float32)} makes every field and every scalar parameter
\code{Float32} (the constructor converts them, and \code{convert\_sliding\_law} converts the law). The code avoids
silent promotion: division guards use \code{eps(T)} rather than a literal \code{1e-15}, \code{accumulate\_psi\_out!} returns
\code{zero(eltype(...))} rather than \code{0.0}, and powers are written \code{x*x} rather than \code{x\^{}2.0}
(\code{Float64\^{}Float64} calls libm's \code{pow}, measured 17 times slower). Three places still compute in
\code{Float64} whatever \code{T} is: the FFT cache buffers, the \code{StaggeredFriction} velocity arrays, and a few
\code{max(0.0, ...)} literals (Section~\ref{pit:open}).

\subsection{Package extensions}\label{arch:ext}
Two weak dependencies (\code{[weakdeps]} in \code{Project.toml}) load code only when the user loads them:
\begin{itemize}[leftmargin=1.5em]
\item \code{CairoMakie} $\to$ \extfile{FastHydrologyMakieExt.jl}: methods for \code{visualize\_field} and
  \code{visualize\_grid}. CairoMakie is slow to compile and not needed to run a model.
\item \code{Shakti} $\to$ \extfile{FastHydrologyShaktiExt.jl}: \code{run!} and \code{step!} for
  \code{TimeSimulation\{<:ShaktiHydroModel\}}. Shakti pulls in CUDA, Metal, AlgebraicMultigrid and more.
\end{itemize}
The functions are declared in \code{src/} without methods (\code{function step! end}, \code{function visualize\_field end}),
so their names belong to FastHydrology and the extension extends them. \code{ShaktiHydroModel\{S\}} is defined in
\file{models/shakti/model.jl} with no reference to Shakti: its field type is a parameter. FastHydrology and Shakti both
export \code{run!} and \code{step!}; a script that loads both must qualify the names, or \code{import} one of them (the tour
uses \code{import Shakti}).

\subsection{Load order}\label{arch:order}
\file{FastHydrology.jl} includes the files in the order of the tour PARTs: the eleven \file{common/} files, then the
seven K24 files, the three HAB files and the Shakti wrapper. The order matters only where a type is named in a struct
field or method signature, which Julia resolves when the definition is evaluated: \file{common/grid.jl} comes first
(nearly every signature names \code{AbstractHydroGrid}); \file{common/state.jl} and \file{common/model.jl} come before
the files whose methods take \code{HydroState} or \code{AbstractHydroModel}; \file{models/kazmierczak2024/model.jl}, which
defines every option type, precedes the files that dispatch on them. Names used only inside function bodies may refer
forward: \code{update\_steady\_state!} in \file{models/kazmierczak2024/run.jl} calls \code{update\_q!} from
\file{water\_flux.jl}, and \code{update\_q\_from\_psi\_out!(model, grid, state)} is called in \file{water\_flux.jl} but defined
in \file{routing.jl}, which comes after it. The tour checks at run time that its PART order is the include order.

\subsection{Dependencies}\label{arch:deps}
\begin{center}\small
\begin{tabular}{>{\raggedright\arraybackslash}p{3.3cm}>{\raggedright\arraybackslash}p{11.8cm}}
\toprule
package & used for \\
\midrule
ImageFiltering & \code{imfilter!}/\code{centered}: the generic \code{convolve!} default (and the tour's reference) \\
FFTW & the cached convolution (\code{plan\_rfft}, \code{plan\_irfft}) \\
OffsetArrays & loaded with ImageFiltering's \code{centered}; no direct use \\
SpecialFunctions & \code{erf} in Eq.~\eqref{eq:N} \\
Statistics, LinearAlgebra & \code{mean} (data loaders), \code{mul!} (FFT) \\
MAT & \code{load\_Kazmierczak} (KORI-ULB \code{.mat} output) \\
NCDatasets & \code{load\_yelmox}, checkpoints, \code{Net\allowbreak{}CDFOutput\allowbreak{}Writer} \\
DocStringExtensions & \code{\$(TYPEDSIGNATURES)} in docstrings \\
\bottomrule
\end{tabular}
\end{center}
Julia $\ge$ 1.10.10 (\code{[compat]}). The test environment (\code{test/Project.toml}) adds MAT, NCDatasets, Shakti (from
GitHub) and pins ParallelStencil to 0.16.0, the version Shakti requires.
